\section{引言与背景回顾}

本节课是 KAIST CS492D ``Design and Learning of Visual Generative Models'' 课程的第五讲，由 Minhyuk Sung 教授主讲。课程的核心主题是 \textbf{Denoising Diffusion Implicit Models (DDIM)}，这是加速扩散模型生成过程的里程碑式工作（Song et al., ICLR 2021）。

\subsection{从 DDPM 到加速采样的动机}

在前几讲中，我们已经学习了 DDPM（Denoising Diffusion Probabilistic Models）的基本框架：通过逐步添加高斯噪声的前向过程（forward process），将数据分布转化为标准高斯分布；再通过训练神经网络学习逆过程（reverse process），从噪声逐步恢复出干净数据。

\begin{importantbox}{DDPM 的核心瓶颈：采样速度}
DDPM 通常需要 $T=1000$ 步的迭代去噪过程来生成一张高质量图像。每一步都需要一次完整的神经网络前向传播，这使得生成速度远慢于 GAN 等单步生成模型。如何在不牺牲质量的前提下减少采样步数，是扩散模型研究的核心挑战之一。
\end{importantbox}

讲者通过对比三类主流生成模型的优劣来引出 DDIM 的研究动机：

\begin{table}[H]
\centering
\caption{主流生成模型特性对比}
\begin{tabular}{lccc}
\toprule
\textbf{模型} & \textbf{样本质量} & \textbf{多样性} & \textbf{生成速度} \\
\midrule
GAN & 高 & 低（模式崩塌） & 快（单步） \\
VAE & 低 & 高 & 快（单步） \\
DDPM & 高 & 高 & 慢（$\sim$1000步） \\
\bottomrule
\end{tabular}
\end{table}

扩散模型在质量和多样性上都表现优异，唯一的短板在于生成速度。DDIM 正是解决这一问题的第一个重要突破。

\subsection{Score-Based Models 与 DDPM 的联系}

讲者在开场部分简要回顾了 score-based generative models 与 DDPM 的关系。两者的核心思想高度相似：

\begin{knowledgebox}{Score Function（得分函数）的核心角色}
对于数据分布 $q(\mathbf{x}_t)$，其 score function 定义为：
$$
\nabla_{\mathbf{x}_t} \log q(\mathbf{x}_t)
$$
在 score-based models 中，我们训练一个神经网络 $s_\theta(\mathbf{x}_t, t)$ 来近似这个 score function。通过 Langevin dynamics 迭代采样，可以从任意先验分布出发，逐步趋近数据分布的高密度区域。
\end{knowledgebox}

Score-based models 的基本流程是：
\begin{enumerate}
    \item 对原始数据分布 $q(\mathbf{x}_0)$ 添加不同程度的高斯噪声，得到一系列模糊分布 $q(\mathbf{x}_t)$
    \item 训练 score predictor 网络，通过最小化 L2 损失来匹配每个噪声水平下的 score function
    \item 采样时从高噪声水平开始，利用 annealed Langevin dynamics 逐步降低噪声水平
\end{enumerate}

\begin{warningbox}{多尺度噪声的必要性}
如果只添加少量噪声（小方差），score function 在远离数据点的低密度区域几乎为零，网络既难以准确学习、也无法提供有效的梯度引导采样点移动。如果只添加大量噪声（大方差），虽然覆盖范围广，但模糊后的分布与原始分布偏差很大。因此，必须使用从大到小的多尺度噪声调度（noise schedule），先粗略定位、再精细恢复——这正是 DDPM 和 score-based models 共享的核心设计思想。
\end{warningbox}

讲者指出，DDPM 的训练目标（噪声预测损失 $\|\boldsymbol{\epsilon} - \boldsymbol{\epsilon}_\theta(\mathbf{x}_t, t)\|^2$）实际上等价于 score matching 的目标。Score function 与噪声预测器之间存在如下关系：

$$
\nabla_{\mathbf{x}_t} \log q(\mathbf{x}_t | \mathbf{x}_0) = -\frac{\boldsymbol{\epsilon}}{\sqrt{1 - \bar{\alpha}_t}}
$$

其中 $\boldsymbol{\epsilon}$ 是添加的噪声，$\bar{\alpha}_t$ 是累积噪声调度参数。因此，噪声预测网络 $\boldsymbol{\epsilon}_\theta$ 本质上就是在学习 score function 的缩放版本。

\subsection{连续时间视角的预告}

讲者还提到，当我们将扩散过程推广到连续时间域时，前向过程可以描述为随机微分方程（SDE）：
$$
d\mathbf{x} = f(\mathbf{x}, t)\,dt + g(t)\,d\mathbf{w}
$$
其逆过程同样有闭式 SDE 表达，且核心仍依赖 score function。DDPM 和 score-based models 都可视为对这个连续过程的不同离散化方案。

\begin{knowledgebox}{统一框架下的不同视角}
DDPM 从概率图模型（马尔可夫链）出发定义前向和逆向过程；score-based models 从 Langevin dynamics 出发学习 score function。两者在数学上等价，仅是同一框架的不同参数化。这一统一视角（Song et al., ICLR 2021, ``Score-Based Generative Modeling through Stochastic Differential Equations''）为后续的 DDIM、DPM-Solver 等加速方法奠定了理论基础。
\end{knowledgebox}

\subsection{本章小结}

本节回顾了 DDPM 与 score-based models 的等价性，明确了扩散模型面临的采样速度瓶颈。DDPM 需要约 1000 步迭代才能生成高质量样本，而本讲将介绍的 DDIM 通过重新定义前向过程（从马尔可夫过程推广到非马尔可夫过程），可以将步数降低到 50 步左右而几乎不损失质量。


